% RESIDUAL-STATE-PARETO1
\section*{RESIDUAL-STATE-PARETO1: promise-aware continuation nonmergeability}

\textbf{STATUS: COMPUTATIONAL finite experiment plus literature-grounded model correspondence. P versus NP remains OPEN.}

\paragraph{Target and model.}
Let \(D(n,8)\) be the exact-weight-eight clause-presence classifier from DN8-FRONTIER1.
Fix a read-once nonadaptive queried set \(Q\) of clause bits. Two legal prefix histories
\(h,h'\) are \emph{continuation-incompatible} if there exists a common legal suffix
continuation \(z\) for which the completed formulas have opposite SAT/UNSAT outputs.
Any deterministic residual-state machine respecting this cut must put incompatible histories
in distinct states.

The exact-weight promise changes the state problem materially. If two histories contain
different numbers \(k\) of selected clauses, then they have disjoint legal continuation
domains and can share a state. Thus explicitly retaining \(k\) as mandatory memory is not a
valid promise-domain lower-bound argument.

For fixed \(k\), let the residual signature of a prefix be the set of suffix selections that
complete it to an UNSAT exact cover. Distinct residual signatures are pairwise
continuation-incompatible. Hence the exact number of live nonterminal states required at
that cut is the maximum, over feasible \(k\), of the number of distinct nonempty residual
signatures. A zero signature is terminal SAT and need not occupy a live internal state.

\paragraph{Literature correspondence.}
Bova--Slivovsky's subterm width for OBDDs is the closest total-function analogue: it counts
distinct residual CNFs/subfunctions induced by prefixes, and their Theorem~4 constructs an
OBDD of size at most the number of variables times subterm width. Their expander
construction also yields exponential OBDD lower bounds from large subfunction width.
Our object differs because the input variables are clause-presence bits and the function is
partial under an exact-Hamming-weight promise.

Dinesh--Koroth--Sarma define bitwise decomposable projective dimension and prove that it
characterizes deterministic branching-program size up to polynomial factors:
\[
\operatorname{bpdim}(f)/6\le \operatorname{bpsize}(f)
\le \operatorname{bpdim}(f)^c
\]
for an absolute constant \(c\). This is a substantially stronger target than ordinary
projective dimension, for which they exhibit exponential gaps to branching-program size.

Beame--Jayram--Saks prove genuine time--space tradeoffs for general branching programs.
Pebbling and cumulative-memory results provide a separate exact abstraction of the
STORE-versus-RECOMPUTE issue. These results motivate a joint time/state/representation
frontier, but none currently transfers the present static-SAT cut certificate to arbitrary
polynomial-time RAM computation.

\paragraph{Exact finite experiment.}
The verifier is
\texttt{experiments/universal-dag/residual\_state\_pareto.py}
(commit \texttt{dd86352b1ab48e0635b351d36d8c0a3560262c08}).
It regenerates the 272 exact covers for \(D(4,8)\) and 4,380 exact covers for
\(D(5,8)\), and then computes exact residual signatures and their
\(\mathbb F_2\)-incidence rank at every cut of the displayed schedules.

For \(D(4,8)\), sign-pattern round-robin has peak 32 active states; grouping clauses by
variable triple has peak 56. A heuristically optimized nested schedule has
\[
\boxed{W_{\rm active}^{\rm tested}(4,8)=20}
\]
at levels 15,16,17,19,20, with active-state area 318 and
\[
\sum_t\log_2 W_t=89.439787863.
\]
Including the shared terminal-SAT class gives 21 cut classes.

For \(D(5,8)\), a heuristically optimized nested 80-query schedule has
\[
\boxed{W_{\rm active}^{\rm tested}(5,8)=123}
\]
at levels 54 and 55, active-state area 4,557, and
\[
\sum_t\log_2 W_t=409.958078146.
\]
At the balanced 40-query cut the exact active count is 72.

These schedules are feasible witnesses only; they are not proven globally optimal over all
read-once schedules.

\paragraph{Characteristic-two rank observation.}
At the peak cuts of the best tested schedules, the residual-completion incidence matrix over
\(\mathbb F_2\) has rank exactly equal to the number of active continuation classes:
\[
D(4,8):\quad \operatorname{rank}_{\mathbb F_2}=20=W_{\rm active},
\]
\[
D(5,8):\quad \operatorname{rank}_{\mathbb F_2}=123=W_{\rm active}.
\]
This is a finite COMPUTATIONAL identity, not a theorem equating ordinary rank,
projective dimension, bitwise decomposable projective dimension, or branching-program size.
It nevertheless supplies a concrete state resource to which the project's multisign/
product-span machinery can be compared.

\paragraph{Hostile controls.}
\begin{itemize}
\item Parity has decision depth \(N\) but exactly two nonterminal residual functions at every
proper cut, so the continuation-state measure assigns only one live bit. This survives the
basic affine/Gaussian-elimination sanity check that raw decision counting fails.
\item DN8-FRONTIER1 gives an \(O(n^6)\) pair-conflict circuit for fixed \(m=8\).
Therefore growth of the read-once state measure cannot imply unrestricted SAT hardness.
\item Rereading and RECOMP remain open escape routes. A general branching program may
reuse variables and merge histories in ways excluded by the current fixed read-once schedule.
\end{itemize}

\paragraph{New narrowed arrow.}
The experimental target can now be stated more sharply:
\[
\boxed{
\text{multisign residual incidence}
\Longrightarrow
\text{promise-compatible general-BP state lower bound}
}
\]
rather than
\[
\text{raw decision multiplicity}\Longrightarrow\text{runtime}.
\]
The most concrete next test is a partial-function analogue of bitwise decomposable
projective dimension, or another continuation-incompatibility parameter that permits
adaptive rereading. Only after such a certificate survives RECOMP should it be connected to
pebbling/cumulative memory.

\begin{thebibliography}{9}
\bibitem{BovaSlivovsky17}
S.~Bova and F.~Slivovsky,
\emph{On Compiling Structured CNFs to OBDDs},
Theory of Computing Systems 61:637--655, 2017.
doi:10.1007/s00224-016-9715-z.

\bibitem{DKS16}
K.~Dinesh, S.~Koroth, and J.~Sarma,
\emph{Characterization and Lower Bounds for Branching Program Size Using Projective Dimension},
FSTTCS 2016, LIPIcs 65:37:1--37:14.
doi:10.4230/LIPIcs.FSTTCS.2016.37.

\bibitem{BJS01}
P.~Beame, T.~S.~Jayram, and M.~Saks,
\emph{Time--Space Tradeoffs for Branching Programs},
J. Comput. Syst. Sci. 63(4):542--572, 2001.
doi:10.1006/jcss.2001.1778.

\bibitem{Nordstrom13}
J.~Nordstr\"om,
\emph{Pebble Games, Proof Complexity, and Time--Space Trade-offs},
Logical Methods in Computer Science 9(3):15, 2013.
doi:10.2168/LMCS-9(3:15)2013.
\end{thebibliography}

% END_OF_NOTE::RESIDUAL_STATE_PARETO1
